% ============================================================
% DRAFT NOTE (setup)
% Contains: Grouped splits, prompt weighting, stratified reporting, metrics,
% the secondary quality-failure validation, and the transfer protocol.
% Written now: The evaluation design per the pivot discussion.
% Pending: Frozen split sizes, metric list finalization, and the
% quality-failure validation thresholds, to be preregistered before results
% are inspected.
% Sources: pivot discussion; prior draft (evaluation discipline)
% ============================================================
\section{Experimental Setup}

\todo{STALE SECTION: written for the divergence objective; to be rewritten for the quality-risk target. Do not review yet.}
\label{sec:setup}

\paragraph{Splits and weighting.}
All evaluation is grouped by prompt: every run and every token row of a
prompt lives in exactly one group, so millions of token rows are never
treated as millions of independent samples. Prompts receive equal total
weight. The corpus is split into 508 development and 256 confirmation prompt
groups; every modeling decision, including architecture selection and the
comparator, was frozen on development data, and the confirmation set was
opened and scored exactly once. \todo{Cite \texttt{protocol\_lock.json} and
\texttt{confirmation\_prefit\_lock.json} for the frozen quarantine.}

\paragraph{Stratified reporting.}
Results are reported separately by task, compressor, ratio, horizon, and
early versus late position; a pooled average cannot rescue a stratum that
fails. Primary metrics are rank correlation and regression error of predicted
against realized divergence at each horizon, with prompt-cluster bootstrap
intervals. \todo{Preregister the exact metric list and the decision rule per
stratum before inspecting held-out results.}

\paragraph{Secondary validation: from divergence to quality.}
Divergence is only a useful hazard signal if it corresponds to outcomes a
user cares about. As a preregistered secondary analysis, we test whether
per-run aggregates of predicted damage separate runs with realized
task-quality failures (wrong or degraded final answers, and catastrophic
failure indicators recorded in the corpus) from clean runs. Quality fields
are used only in this validation, never as supervision for the forecaster.
\todo{Freeze the aggregation statistic and separation metric.}

\paragraph{Transfer protocol.}
The forecaster is trained on the primary corpus and evaluated on the
second-family traces in two modes: zero-shot, and after lightweight
adaptation on a small grouped subset. \todo{Specify the adaptation budget and
report both modes once the transfer traces are collected.}
